\chapter{Conclusion and Future work}
\label{chap:conclusion_future_work}

% ------------------------------------------------------------
\section{Conclusions}
\label{sec:conclusions}
% ------------------------------------------------------------

Field cancerisation represents one of the most compelling and, at the same
time, methodologically challenging phenomena in cancer epigenomics.
Histologically indistinguishable from healthy tissue, tumour-adjacent breast
parenchyma harbours early epigenetic alterations that precede overt
neoplastic transformation and whose systematic characterisation may
ultimately inform both risk stratification and early detection strategies.
The central biological question motivating this thesis --- whether DNA
methylation differences between Normal and Normal-Adjacent tissue are
detectable, reproducible, and transferable across independent cohorts under
strict statistical control --- was addressed through an end-to-end
computational pipeline integrating multi-cohort data harmonisation,
biologically informed feature selection, and constrained combinatorial
optimisation.

The study was built on three publicly available DNA methylation datasets
(\texttt{GSE69914}, \texttt{GSE225845}, \texttt{GSE287331}) spanning the
Normal\,--\,Adjacent\,--\,Tumour axis in breast tissue across two Illumina
array generations (450K and EPIC).
Exploratory analysis confirmed that, at the whole-methylome level, Normal
and Adjacent tissues are globally similar, with nearly superimposed
$\beta$-value distributions and comparable variance profiles.
Yet locus-specific analyses consistently revealed a focal early drift in
Adjacent samples: outlier burden ratios ranging from~1.35 in
\texttt{GSE69914} to~13.33 in \texttt{GSE287331}, and silhouette scores up
to~0.34, provided quantitative evidence for the presence of field effects
across all three independent cohorts.
This cross-dataset consistency confirms that the Normal\,--\,Adjacent signal
is not an artefact of a single study design but a reproducible biological
phenomenon, coherent with contemporary models of tumour initiation in which
locally restricted epigenetic perturbations accumulate in histologically
normal tissue prior to overt malignant transformation.

Building on this exploratory foundation, the preprocessing pipeline
addressed the principal sources of confounding in multi-cohort methylation
studies: probe-design bias,
batch effects between the datasets (corrected via NeuroCombat with age as a
continuous covariate, estimated strictly on the training set to prevent
leakage), and age-dependent methylation structure (quantified through the
Horvath multi-tissue epigenetic clock and subsequently controlled via age
binning and CpG blacklisting).
The transformation of raw $\beta$-values into absolute deviations from the
Normal reference $\lvert\beta - \mu_{\mathrm{Normal}}\rvert$ provided a
biologically grounded input space that resolves the drift polarity
heterogeneity identified across cohorts while aligning naturally with the
field cancerisation hypothesis.

Feature selection proceeded through a multi-stage stability-based pipeline:
empirical $\beta$-filtering to remove constitutively non-variable CpGs,
variance filtering in $M$-space, repeated stratified $k$-fold stability
ranking with directional weighting (RSKF, Spearman
$\rho = -0.848$ between score and $|\Delta M|$), region anchoring to CpG
islands, correlation-based redundancy pruning ($|r| \geq 0.85$), and
genomic diversification with chromosomal and spatial constraints.
The output of this pipeline --- 5\,000 stable, non-redundant, genomically
diverse CpG sites --- constituted the input space for a LinearSVC
classifier, whose coefficients
were then passed as relevance scores to the main methodological contribution
of this work.

The MILP knapsack formulation represents the core novelty of this thesis.
Rather than selecting CpG sites by greedily maximising predictive
performance, the optimiser explicitly encodes biological prior knowledge as
hard and soft constraints: enrichment in COSMIC breast cancer genes
($W_{\mathrm{COSMIC}}$), upper bounds on chromosomal representation
(\textsc{chr\_max}), spatial diversity (max per 500\,kb window), redundancy
control (\textsc{corr\_thr}), gene-level diversity (\textsc{max\_per\_gene}),
and a minimum fraction of hypermethylated sites (\textsc{min\_hyper}).
Applied to the multi-cohort pooled setting of Chapter~7, the resulting panel
of 50 CpG sites achieves $\mathrm{AUC} = 0.902$ on the pooled held-out test
set ($n = 177$: 104 Normal, 73 Adjacent), with per-cohort breakdown
$\mathrm{AUC}_{\texttt{GSE287331}} = 1.000$,
$\mathrm{AUC}_{\texttt{GSE225845}} = 0.827$,
$\mathrm{AUC}_{\texttt{GSE69914}} = 0.487$,
against the 5\,000-CpG upper bound of $\mathrm{AUC} = 0.944$.
Of the 50 selected CpGs, 18 overlap with the COSMIC Cancer Gene Census
breast-cancer subset, compared to~11 in the unconstrained baseline
($\mu = 0$), demonstrating the effectiveness of the knapsack biological
penalty.

The biological interpretation of the panel corroborated the coherence of
the constrained selection.
The 50 CpGs map to 49 unique genes spanning tumour suppressors
(\textit{RB1}, \textit{TP53}), chromatin regulators (\textit{HDAC4},
\textit{NCOR1}), developmental transcription factors (\textit{TBX3},
\textit{SALL4}, \textit{GATA3}), signalling components (\textit{NOTCH1},
\textit{NTRK3}, \textit{AKT1}), DNA repair (\textit{BRCA2}), and an
RNA-binding oncogene (\textit{LIN28B}) --- a multifactorial portrait of
epigenetic field reprogramming that could not have emerged from a purely
performance-driven selection.
This is the central result of the thesis: \emph{interpretability is not a
by-product of performance --- it must be structurally imposed}.

Taken together, these results demonstrate that it is possible to construct a
compact, clinically motivated CpG panel that preserves near-baseline
classification performance while embedding explicit biological knowledge.
The MILP knapsack framework is fully parametrisable: the constraint
structure can be adapted to any tumour type or epigenetic phenotype for
which a reference cancer gene database exists, making the approach general
beyond the breast cancer context explored here.



% ------------------------------------------------------------
\section{Future Work}
\label{sec:future_work}
% ------------------------------------------------------------

Several directions naturally extend the present work, ranging from
immediate experimental validations to broader methodological
generalisations.

\paragraph{External validation on independent cohorts}
The most pressing limitation of this study is the absence of evaluation on
a dataset not involved in any step of the pipeline design.
The three cohorts employed here were all used in training and feature
selection, albeit with strict leakage controls.
Identifying an additional dataset with Normal and Adjacent breast
tissue samples --- a non-trivial task given the scarcity of publicly
available adjacent-normal methylation data with consistent phenotype
annotation --- would provide a genuine out-of-distribution test of the
50-CpG panel.
Cross-platform validation, extending the panel to the EPIC~v2 array
(${\sim}930$k CpG, broader promoter coverage), would further establish its
robustness across array generations.

\paragraph{Sensitivity analysis on pipeline and MILP hyperparameters}
A systematic sensitivity analysis of the feature selection pipeline
hyperparameters --- including the directional stability threshold, the
biological prior weight $\lambda$, the correlation redundancy cutoff
$\tau_c$, and the diversification target $K$ --- would provide empirical
grounding for the design choices adopted in Chapter~5.
Given the computational cost of re-running stability selection across the
full hyperparameter grid, this analysis would benefit from a reduced
resampling scheme or surrogate evaluation on a representative subset of
CpGs.
Analogously, varying each MILP constraint independently ($\mu$,
$W_{\mathrm{COSMIC}}$, \textsc{chr\_max}, \textsc{corr\_thr},
\textsc{max\_per\_gene}, \textsc{min\_hyper}) while holding the others
fixed would quantify which constraints drive biological enrichment and which
are partially redundant, providing practical guidance for adapting the
framework to other cancer types.

\paragraph{Replacement of SVM coefficients with transformer attention
weights (CpGPT)}
In the present formulation, the relevance scores fed to the knapsack are
the linear coefficients of a LinearSVC trained on the 5\,000 candidate
CpGs.
While interpretable and computationally efficient, linear coefficients
capture only additive, univariate contributions and may miss higher-order
co-methylation interactions.
A natural extension is to replace the SVM with a transformer-based model
pre-trained on large-scale methylation data, such as
CpGPT~\cite{levy2024cpgpt}, and to use the resulting attention weights as
the relevance scores for the MILP objective.
The self-attention mechanism~\cite{vaswani2017attention} learns
context-dependent representations where the contribution of each CpG is
modulated by its co-methylation context across the genome; feeding these
richer relevance estimates into the knapsack would preserve the full
biological constraint structure while potentially improving both
discriminative power and biological coherence of the selected panel.

\paragraph{Multi-omic integration}
DNA methylation captures only one layer of the epigenetic regulation
underlying field cancerisation.
The CpG sites identified by the knapsack panel represent natural candidates
for functional validation through complementary omics modalities.
Integration with RNA-seq data from matched Normal\,--\,Adjacent\,--\,Tumour
samples would allow assessment of whether hypermethylation at the selected
CpGs is associated with transcriptional silencing of the corresponding
genes, directly linking the epigenetic signal to downstream expression
changes.
Similarly, ATAC-seq or ChIP-seq data for histone marks (H3K27me3,
H3K4me3) at the selected loci would clarify whether the observed
methylation alterations co-occur with chromatin accessibility changes,
providing a more complete picture of the epigenetic state of adjacent
tissue.

\paragraph{Generalisation to other tumour types and phenotypes}
The MILP constraint structure is fully parametrisable.
Applying the same pipeline to colorectal, lung, or ovarian cancer ---
tumour types for which large EPIC methylation datasets with normal and
adjacent tissue are increasingly available --- would test whether
biologically constrained CpG panel selection is a generally effective
strategy or whether its advantages are specific to the breast cancer
context.
More broadly, the framework could be adapted to phenotypes other than
tissue state classification: age-acceleration groups, molecular subtypes,
or treatment response endpoints could serve as the target variable, with
the COSMIC constraint replaced by phenotype-relevant gene sets from MSigDB
or other curated databases.

% ------------------------------------------------------------
\section{Code and Data Availability}
\label{sec:availability}
% ------------------------------------------------------------

All code developed for this thesis, including preprocessing pipelines,
feature selection notebooks, MILP formulation, and evaluation scripts, is
publicly available at:
\begin{center}
    \url{https://github.com/elisabettaroviera/THESIS}.
\end{center}


%A systematic sensitivity analysis of the feature selection 
%pipeline hyperparameters — including the directional 
%stability threshold, the biological prior weight $\lambda$, 
%the correlation redundancy cutoff $\tau_c$, and the 
%diversification target $K$ — would provide empirical 
%grounding for the design choices adopted in 
%Chapter~\ref{chap:feature_selection}. Given the 
%computational cost of re-running stability selection 
%across the full hyperparameter grid, this analysis 
%would benefit from a reduced resampling scheme or 
%surrogate evaluation on a representative subset of CpGs.